\section{The electric contour integral}\label{sec:04-gas-fluctuations}

The quadratic integration estimate is real. Electric sectors also carry
an oscillatory imbalance factor, and obtaining its sharp contribution
requires a complex displacement. We first implement the deterministic
trial fields by moving the actual particles, while keeping their integer
dimensions fixed. The next section then varies the dimensions through
an explicitly constructed holomorphic approximation.


\subsection{Reference dimensions and compensated quadratic costs}

Freeze the positions, labels, and counts of the explicit particles
\(\zeta\) from Section~\ref{sec:03-gas-mean}. For a pair of
nonnegative integers \(\nu=(\nu_+,\nu_-)\) and an independent
parameter \(M\in\C\), define
\[
 F_\nu(M)=\int A(X_\nu,\zeta)\exp(-\Phi_M(q)),
 \qquad q_\sigma=X_{\nu,\sigma}-\nu_\sigma p_L,
\]
with the two unnormalized reference measures from
Lemma~\ref{lem:03-compensated-sector}. Thus \(F_\nu\)
includes their masses \(J_{\nu_+}(w_0)J_{\nu_-}(w_0)\), but not
the extracted site product, the explicit-particle prefactor, or the
imbalance phase. The energy \(\Phi_M\) retains the prescribed
trial totals \(m'(M)\); the empirical field is centered at the
actual integer dimensions \(\nu\). These two quantities need not
agree. Their independence will be needed for interpolation in the
next section.

Write
\[
 d_e=\ell_e+\frac{(\operatorname{Re}M)^2+n_e^2}{L},
 \qquad H_0=H_{\rm ref}.
\]
We use \(H_0\) in the \(s\)-coordinate on both groups.
Derivatives of primitives are densities or distributions in that
coordinate and are pushed forward by \(W\) when inserted in
\(\mathcal E\).

\begin{proposition}[Bounds at integer reference dimensions]\label{prop:04-integer-bounds}
Fix \( |\operatorname{Re}\eta|\le a_0<\pi\) and
\( |\operatorname{Im}\eta|\le b_0<\infty\), with the parameter
conventions of Proposition~\ref{prop:03-trial-bounds}. For
nonnegative integer dimensions \(\nu_\sigma\), put
\[
 \mathcal L_{\nu,j}(M)
 =\log|e^{i(\eta+\pi f+2\pi j)M}F_\nu(M)|,
 \qquad v=\operatorname{Im}M.
\]
The following estimates hold.
\begin{enumerate}
\item If \(M\) is real and \(0\le\nu_\sigma\le C e^{\gamma L}\),
then, without any restriction on the explicit counts,
\[
 \mathcal L_{\nu,j}(M)\le C_hL+C_hn_e-c_hd_e.
\]
This bound is independent of the integer alias \(j\).
\item Suppose \(n_e\) is polynomially bounded, as in part 2 of
Proposition~\ref{prop:03-trial-bounds}. In the same dimension range,
for every real \(v\) and integer \(j\),
\begin{equation}
 \mathcal L_{\nu,j}(M)\le C_h L+C_h n_e-c_h d_e-2\pi jv
                 +C_h(v^2+(1+f)|v|).
 \label{eq:G-15}
\end{equation}
\item Assume the common hypotheses of parts 3 and 4 of
Proposition~\ref{prop:03-trial-bounds}, polynomial bounds on
\(n_e,|j|,|M|,\ell_e\), and \(\nu_\sigma\asymp m\).
Choose the imaginary imbalance and trial density from part 3 of
that proposition. Then
\begin{equation}
 \mathcal L_{\nu,j}(M)\le C_h L+C_h n_e-c_h d_e-c_\delta U^2 L+o(L).
 \label{eq:G-16}
\end{equation}
If \(U\le U_0\), the choice from part 4 instead gives
\begin{equation}
 \mathcal L_{\nu,j}(M)\le\frac{3\gamma L}{8}
 -2L\frac{\operatorname{Re}(\eta^2)}{4\kappa}
 +C\varepsilon L+o(L)+C_h n_e-c_{h,\varepsilon}d_e.
 \label{eq:G-17}
\end{equation}
Here \(\delta\) can be taken sufficiently small without degrading
the displayed damping and per-particle constants, and \(U_0\) may
then be fixed arbitrarily, as in part 4 of that proposition.
\end{enumerate}
The estimates are uniform once the parameters in the trial bounds
and the polynomial size powers are fixed.
\end{proposition}

We prove Proposition~\ref{prop:04-integer-bounds} using the real integration bound of Lemma~\ref{lem:04-quadratic} and the particle deformation constructed below.

Throughout this section, a polynomial bound means a bound by a fixed power of \(L\), whose exponent may depend on previously fixed parameters. It never refers to the exponentially large reference dimensions. Energy expressions use the complex bilinear extension of \(h_0\), whereas norms and adjoints use the Hermitian convention. Angle brackets against physical test functions denote distributional pairing.

For smooth real Hilbert fields \(Y\), and a real smooth primitive \(Z_0\) compactly supported in the core \(|s|\le(1-\delta_c)L\) for some fixed \(\delta_c>0\), use the functional
\[
 C_z(Y)=\mathcal E\big(m'p_L+(Y+iZ_0)',\zeta\big)-i(\eta+\pi f+2\pi j)M
       =c_z+h_0(u_z,Y)+h_0(Y,TY)/2
\]
with complex bilinear convention; write \(C_0\) for \(Z_0=0\), so \(C_z(Y)=C_0(Y+iZ_0)\). In \eqref{eq:G-16},\eqref{eq:G-17}, the pushforward of \(Z_0'\)
by \(W\) is the density \(z_\sigma+\sigma v p_L\) chosen in
parts 3 and 4 of Proposition~\ref{prop:03-trial-bounds}. Its required fixed-order derivative bounds are polynomial. Everywhere in the polynomial regimes \(u_z\) is a bounded Hilbert vector of polynomial norm, obtained by duality from a smooth odd-reflected test in \(s\) with polynomial fixed-order derivatives. Indeed use Lemma~\ref{lem:03-source-regularity}, and the action of \(H_{\rm form}\) on smooth core fields (localized diagonal log convolved by their smooth derivatives, as in the core-subspace proof). Also \(c_z\) is polynomially bounded except for the unchanged explicit-explicit constant \((\zeta,H\zeta)_{\rm off}/2\), whose real part is at least \(-{\rm poly}(L)\); the equality in polynomial size for \(c_z\) here uses \(|j|\) polynomial as well. Without polynomial truncation we will only use \(Z_0=0\); the forms are still continuous at fixed sizes and generic frozen positions. The infimum
\[
 m_z=\inf_{Y\ {\rm real\ in}\ \mathcal X}\operatorname{Re} C_z(Y)
\]
satisfies the indicated trial lower bounds (with \(B(x)\) replaced using \eqref{eq:G-3} by positive damping in \(d_e\), minus a constant). Indeed apply them first on dense smooth Dirichlet fields, then use continuity. Every such variation preserves the two specified totals. Primitives correspond to \(s\)-densities here, so integrable square-root singularities at physical endpoints do not impede the energy bounds.

The real reference bound already handles parts 1 and 2. For the sharp
estimates in part 3, the imaginary trial density and the relative
determinant must both be realized inside the original particle integral.
The following deformation does this while leaving its integer dimensions
unchanged.

\begin{lemma}[Particle contour deformation]\label{lem:04-contour}\label{E:particle-transport}
Under the hypotheses of part 3 of
Proposition~\ref{prop:04-integer-bounds},
choose the core projection \(P\) of \eqref{eq:G-10}, with
\(S=P\mathcal X\), \(R=S^\perp\), and Schur complement \(T_e\).
For every prescribed fixed tail exponent, a sufficiently large fixed
\(D\) gives an exact contour deformation of the reference integral.
On \(\|Q\|_{-2,A}\le L^D\), its compensated cost is
\(C_0(Q+a(Q))-h_0(Q,Q)/2+o(1)\), where
\[
 Q+a(Q)=iZ_0+Q_R+T_{SS}^{-1/2}Q_S
             -T_{SS}^{-1}Pu_z-T_{SS}^{-1}T_{SR}Q_R.
\]
The extra Jacobian, after the one-coordinate reference-density
cancellation, has modulus
\((1+o(1))|\det T_{SS}|^{-1/2}\). The full cost \(\operatorname{Re}C_0(Q+a(Q))\) has Hessian
\(I\) on \(S\) and \(\operatorname{Re}T_e\) on \(R\), and its
infimum is at least \(m_z\). The complementary contribution is smaller
than the prescribed stretched exponential.
\end{lemma}

\begin{proof}
We first construct an affine path of smooth fields with a positive real
Hessian. We then implement this path by particle displacements, compute
the compensated cost and the remaining Jacobian, and estimate the region
where a cutoff switches off the deformation.

\paragraph{The affine path and its real cost.}
Choose \(P,S,R,T_e\) as in \eqref{eq:G-10},\eqref{eq:G-11} with \(\xi\) sufficiently small (here \(R=S^\perp\)). After reducing \(\delta_c>0\), all core basis and shift functions are supported in \(|s|\le(1-\delta_c)L\). The frozen affine shifts \(Q\mapsto Q+a(Q)\) along the successive stages are:
\begin{enumerate}
\item Add \(i\tau Z_0\), \(0\le\tau\le1\).
\item Keeping \(i Z_0\) added, replace \(Q_S=P Q\) by \(Q_S-\tau(T_{SS}^{-1}P u_z+B_s Q_R)\), where \(B_s=T_{SS}^{-1}T_{SR}\) here denotes a block map, \(Q_R=(I-P)Q\).
\item In the fully shifted coordinates of stage 2 replace the free input \(Q_S\) of that expression by \(D_\tau Q_S\), with
\[
 D_\tau=((1-\tau)I+\tau T_{SS})^{-1/2}.
\]
\end{enumerate}
Thus the output primitive in the last stage equals \(i Z_0+Q_R+D_\tau Q_S-T_{SS}^{-1}P u_z-B_s Q_R\).
Use a single smooth path parameter in \([0,1]\) by flattening at stage boundaries. Matrix functions and their required derivatives here have polynomial norms by \eqref{eq:G-8}. For instance the symmetric inverse square root is given by the positive-resolvent integral \(\pi^{-1}\int_0^\infty y^{-1/2}(yI+(1-\tau)I+\tau T_{SS})^{-1}dy\), using accretivity (its identity follows from the scalar formula and derivatives on Jordan blocks). All nontrivial input uses of \(Q\) in \(a\) are through smooth dual tests with polynomial bounds.

The real Hessian of the full cost \(\operatorname{Re}C_0(Q+a(Q))\)
has an inverse-polynomial gap throughout this path. To check this
without treating complex changes as real orthogonal maps, write
\(A_S=T_{SS}\), \(C_S=T_{SR}\), and \(B_s=A_S^{-1}C_S\).
The linear part of stage 2 sends \((s,r)\) to \((s-\tau B_sr,r)\).
Direct expansion of the complex bilinear quadratic form gives
\[
 T[(s-\tau B_sr,r),(s-\tau B_sr,r)]
 =T[(s,(1-\tau)r),(s,(1-\tau)r)]
       +(2\tau-\tau^2)T_e[r,r].
\]
For real \(s,r\), both terms have positive real part, and the
coefficients on \(r\) add to one. Stage 1 leaves the Hessian unchanged.
At stage 3 it is block diagonal, with blocks \(T_e\) and
\(A_SH_\tau^{-1}\), where \(H_\tau=(1-\tau)I+\tau A_S\).
The identity
\[
 H_\tau^*\operatorname{Re}(A_SH_\tau^{-1})H_\tau
       =(1-\tau)\operatorname{Re}A_S+\tau A_S^*A_S
\]
and the polynomial norms and inverse norms of these matrices give the
required gap. Here \(D_\tau\) commutes with \(A_S\) and is symmetric,
so \(D_\tau^{\mathsf t}A_SD_\tau=A_SH_\tau^{-1}\).

All residual kernels, linear dual tests, and their path derivatives
have polynomial physical derivative bounds, by the core-subspace
construction. Their trace norms are polynomial as well, since the
change from \(T-I\) has polynomial norm and rank \(O(L)\). At the
endpoint the Hessian is \(I\oplus\operatorname{Re}T_e\).
For the infimum comparison, fix a real \(r\). The stationary point of
\(C_z(s+r)\) in the complex \(S\)-variable is
\(s_*=s_1+is_2=-A_S^{-1}(Pu_z+C_Sr)\). Completing the complex square
at \(s_*\) yields
\[
 \operatorname{Re}C_z(s_*+r)
 =\operatorname{Re}C_z(s_1+r)
       +\tfrac12h_0(s_2,\operatorname{Re}A_Ss_2)
 \ge m_z.
\]
The remaining endpoint input contributes \(\|s\|^2/2\). This proves
the infimum claim even though the final contour need not have the
initial imaginary trial profile.

\paragraph{Realizing the path as an exact contour deformation.}
Choose a \(C^\infty\) function \(\chi:[0,\infty)\to[0,1]\), equal
to one on \([0,1]\) and zero on \([4,\infty)\). If the flattened
path just constructed is \(a_\tau\), use the particle-dependent
parameter
\[
 \tau(Q)=\chi(\|Q\|_{-2,A}^2/L^{2D}),
 \qquad a(Q)=a_{\tau(Q)}(Q),
\]
and homotope from zero by replacing \(\tau(Q)\) with
\(u\tau(Q)\), \(0\le u\le1\). The empirical primitive, as a function
of a particle at \(s_i\), has first derivative \(-\delta_{s_i}\)
and second derivative \(\delta'_{s_i}\). Both belong continuously
to the physical \(H^{-2}\) space: in the sine basis their squared
norms are bounded respectively by constant multiples of
\(A^{-1}\sum_k(1+(k/A)^2)^{-2}\) and
\(A^{-1}\sum_k(k/A)^2(1+(k/A)^2)^{-2}\).
Thus the deformation is \(C^2\) in the real particle variables,
which is sufficient for Stokes' theorem.

Move a particle of label \(\sigma\) at \(s\) to
\[
 \Psi_\sigma(s;Q)=s+\Delta_\sigma(s;Q),\qquad
 \Delta_\sigma(s;Q)=-\frac{a_\sigma(Q)(s)}
                              {\nu_\sigma\lambda_L(s)}.
\]
For any fixed derivative order, on the support of this motion,
\[
 \|\partial_s^k\Delta\|_\infty
 \le L^{C_{D,k}}e^{-\gamma\delta_c L}.
\]
Indeed \(a\) and its required derivatives are polynomial on
\(\|Q\|_{-2,A}\le2L^D\), and
\(\nu_\sigma\lambda_L(s)\ge c e^{\gamma\delta_cL}\) in its fixed
core. This estimate concerns spatial differentiation at frozen \(Q\).

For completeness, the exact change of contour is a change of an
integration chain, not an assertion that the map \(\Psi\) is
holomorphic or globally injective. Partition every original real
coordinate into a slightly enlarged core interval and its complement,
and fix the complementary coordinates. The factor \(A\) is independent
of all active core coordinates. The remaining density times the
coordinate volume form is a holomorphic top form \(\Omega\) in a
neighborhood of the active intervals and their images. Its reference
part is a squared Vandermonde times the holomorphic one-body density,
so coincidences are zeros rather than poles. The regular exponent is
holomorphic there by the core regularity estimates of the preceding
section. Hence \(d\Omega=0\). On a side face of the homotopy prism,
the corresponding coordinate lies at an interval endpoint and its
displacement vanishes in a neighborhood; its differential is zero
on that face, so the pullback of \(\Omega\) vanishes. Stokes' theorem
therefore equates the two integrals. Apply this argument for almost
every frozen configuration and sum the finitely many partitions.
The estimates below give absolute integrability; configurations with
coincident explicit particles may be omitted as null sets.

\paragraph{The compensated cost after displacement.}
The contour identity is exact. We now estimate its pullback uniformly on \(\|Q\|_{-2,A}\le 2L^D\). Write \(X=X_\nu\), \(\Psi={\rm Id}+\Delta\). Separate the fixed-field single-coordinate Jacobians \(1+\Delta_\sigma'\). The diagonal omission in the Vandermonde log change contributes half the diagonal of
\[
 H_0(\Psi(s),\Psi(v))-H_0(s,v)
\]
per particle (kernels here in \(s\), with appropriate labels). Its exponential cancels the transformed \(2\sin\theta\,\theta'\) ratio and these single-coordinate Jacobians exactly. The difference kernel is defined by the log of the cos-difference ratio, with analytic branch near 1, and its required derivatives are exponentially small with odd/even reflections as appropriate. Indeed near a moved diagonal use the divided-difference bounds in the analytic core and small derivatives of \(\Delta\); off diagonal use the sine separation bounds with displacements relative to \(\theta'\). Thus the resulting base negative-log change is half the difference kernel on \(X,X\). Combining with \(\Phi_M\) evaluated after movement and including the phase in the cost, one obtains uniformly up to \(o(1)\)
\begin{equation}
 C_0(Q+a(Q))-\tfrac12 h_0(Q,Q).
 \label{eq:G-19}
\end{equation}
This notation always means the compensated smooth expansion, not the two singular forms separately on an empirical field.

We justify the error in \eqref{eq:G-19} by a single scalar Taylor
expansion, which avoids any dimension-dependent multivariable
remainder. Freeze \(Q\) and \(a\), and replace \(\Delta\) by
\(z\Delta\). On the disk
\( |z|\le R_L=\exp(\gamma\delta_cL/4)\), all required spatial
derivatives of \(z\Delta\) remain exponentially small. Core
analyticity and the divided-difference estimates keep every logarithmic
ratio in its branch near one, including at coincident particle
positions. After removing the unchanged explicit--explicit constant,
the logarithmic cost on this disk is bounded by
\(L^C e^{2\gamma L}\): every term is linear or quadratic in measures
of total variation \(O(e^{\gamma L})\), with polynomially bounded
kernels on the disk. Cauchy's estimate bounds the remainder after
order \(K\), at \(z=1\), by
\[
 C L^C\exp\!\left(\left[2\gamma-
              (K+1)\gamma\delta_c/4\right]L\right).
\]
Choose the fixed integer \(K\) so that \(K+1>8/\delta_c\).

Each retained coefficient is a polynomial in the components of the
frozen displacement in the real span of the core basis and \(Z_0\).
It is therefore enough to compute it for real displacements and extend
the resulting polynomial identity to complex coefficients. For real
displacements, the pushforward of the smooth mean has coefficients
\[
 d_k=\frac{(-\partial_s)^k}{k!}
               (\nu_\sigma\lambda_L\Delta_\sigma^k).
\]
Here \(d_1=a_\sigma'\), while all fixed-order derivatives of \(d_k\),
\(k\ge2\), are exponentially small. The useful cancellation is
\(\nu_\sigma\lambda_L\Delta_\sigma^k=-a_\sigma\Delta_\sigma^{k-1}\);
using the total variation of the mean alone would not prove this
assertion near the edge of the core. All these densities are supported
in the core. For the centered distribution \(q=Q'\), every
positive-order push coefficient, tested against a smooth \(f\), is
\(\langle q,\Delta^k f^{(k)}\rangle/k!\). Integration by parts and
duality with \(\|Q\|_{-2,A}\) make it exponentially small, with all
fixed test derivatives needed below.

These rules give the substitution \(q\mapsto q+a'\) directly for the
smooth kernel \(K\) and the smooth linear term in \(\Phi_M\).
For the base-kernel difference, its \(q,q\) term is exponentially
small by the ratio-kernel bounds and Sobolev duality in both variables.
Terms containing means use \(H_0(\lambda_Lds)=0\). The one-mean
coefficients reduce to a pushed \(q\) paired with \(H_0d_k\),
and the two-mean coefficients pair positive-order \(d_k,d_l\).
These identities first hold for real diffeomorphisms. Differentiation
is legitimate because the changed mean has smooth core density;
local convolution with the diagonal logarithm transfers derivatives
to that density. It gives polynomial smooth norms for \(H_0d_k\).
Consequently the base half-difference is
\[
 (q,H_0a')+\tfrac12(a',H_0a')+o(1).
\]
Polynomial continuation of the finite jets, followed by the scalar
remainder estimate, proves the same identity for complex displacements.
No pushforward density on the real line is asserted for a complex
contour.

Explicitly \(\Phi_M(q)\) is a constant plus its smooth linear test on \(q\) plus \((q,K q)/2\). Hence replacing \(q\) by \(q+a'\) there (and subtracting the phase log) gives exactly \eqref{eq:G-19} minus \((q,H_0 a')+(a',H_0 a')/2\), without ever needing a bare self energy of \(q\). All the Taylor estimates are uniform even close to particle coincidences: on the growing scalar disk the shifted cos-difference over the original one in a same-label base ratio remains near 1, by the simultaneous divided difference (the same frozen shift function is used on the two arguments and has tiny spatial derivatives). The huge constant from a possible very close explicit-explicit pair never participates in these jets.

\paragraph{The Jacobian from the dependence on the configuration.}
The preceding cancellation includes each single-coordinate Jacobian at
fixed \(Q\). It remains to estimate the part arising from
the dependence of \(a\) on the whole configuration. On the full-shift
region write
\[
 D_Qa=\sum_{b=1}^{r}e_b\otimes\phi_b,\qquad r=\dim S=O(L),
\]
where \(\phi_b\) are smooth physical dual tests, and put
\(J_i=1+\partial_s\Delta_{\sigma_i}(s_i;Q)\). The sign
\(\partial_{s_i}Q=-\delta_{s_i}\) cancels the minus sign in
\(\Delta=-a/(\nu\lambda_L)\). The matrix determinant lemma gives
exactly, after factoring \(\prod_iJ_i\), the reduced determinant
\[
 \det(I+V),\qquad
 V_{bb'}=\sum_i
 \frac{\phi_{b,\sigma_i}(s_i)e_{b',\sigma_i}(s_i)}
      {J_i\nu_{\sigma_i}\lambda_L(s_i)}.
\]
For the summand test \(g=\phi_b e_{b'}/(J\nu\lambda_L)\), centered
quadrature is quantified by
\[
 \left|\sum_{i:\sigma_i=\sigma}g(s_i)
                     -\nu_\sigma\int\lambda_Lg\,ds\right|
 =|\langle Q_\sigma,g'\rangle|
 \le \|Q_\sigma\|_{-2,A}\|g'\|_{H^2}
 \le L^C e^{-cL}.
\]
The support of \(e_{b'}\) is inside the core, so all norms on the
right have the required exponential factor. Replacing \(J\) by one
in the mean term has the same error. Thus
\(V_{bb'}=\int\phi_b e_{b'}\,ds+O(L^Ce^{-cL})\), with the
labels summed. On \(S\), the derivative of the endpoint affine map
is \(D_1\); hence \(I+V=D_1+E_L\),
\(\|E_L\|\le L^Ce^{-cL}\). Since \(\|D_1^{-1}\|\) is polynomial,
\[
 \big|\log\det(I+D_1^{-1}E_L)\big|
 \le 2r\|D_1^{-1}E_L\|=o(1).
\]
The principal matrix square root satisfies
\(|\det D_1|=|\det T_{SS}|^{-1/2}\), giving the asserted modulus.

\paragraph{The transition region and the outer tail.}
In the transition region, differentiation of the cutoff adds a single
rank-one term with output \(\partial_\tau a_\tau\in S+\C Z_0\).
The total rank is at most \(r+1\). The cutoff differentials on point
deltas have polynomial bounds on this region, and every entry of the
reduced matrix is at most \(e^{C_DL}\), even if bounded without
centering. Hadamard's inequality therefore bounds its modulus by
\(\exp(O_D(L^2))\). This estimate uses the reduced matrix of size
\(O(L)\), never the exponentially large particle matrix. Outside
the cutoff region the extra determinant is one.

Finally fix a path parameter \(\tau\). Its compensated real cost has
inverse-polynomial gap and polynomial linear and trace-norm bounds,
with constant bounded below by minus a polynomial. These bounds are
independent of the cutoff exponent \(D\). Choose \(E\) once so large
that subtracting \(L^{-E}\|Q\|_{-2,A}^2\) preserves half the gap.
Lemma~\ref{lem:04-quadratic} and the elementary inequality
\(\mathbf1_{\{\|Q\|>L^D\}}\le
 \exp(L^{-E}\|Q\|^2-L^{2D-E})\) yield
\[
 \E_0[\mathbf1_{\{\|Q\|_{-2,A}>L^D\}}e^{-V_\tau(Q)}]
 \le \exp(-L^{2D-E}+L^{E_1}),
\]
with \(E_1\) fixed independently of \(D\). For example, the
negative log sum is at most a polynomial because
\(-\log\lambda\le g^{-1}(1-\lambda)\) on \([g,1]\), while the
trace norm and inverse gap are polynomial.

The parameter in the transition region depends on \(Q\), so a
fixed-parameter estimate alone is insufficient. On
\(\|Q\|_{-2,A}\le2L^D\), the physical kernel bounds give
\(|\partial_\tau V_\tau(Q)|\le L^{C+2D}\), after factoring out
any unchanged explicit-pair constant. A parameter grid of mesh
\(L^{-C-2D-1}\) approximates every such cost within one. Bound its
exponential by \(e\) times the sum over grid costs and apply the
preceding tail estimate at every grid point. The additional logarithmic
cost is only \(O_D(\log L)\). The undeformed outer tail already
has fixed parameter zero. Combining this estimate with the transition
Jacobian bound and bounded reference masses gives
\[
 \exp\{-L^{2D-E}+L^{E_1}+O_D(L^2)\}.
\]
Increasing \(D\) makes this smaller than any prescribed
\(\exp(-L^{E_2})\). The factor \(A\), which depends only on fixed
coordinates, remains in \([0,1]\) in all these absolute estimates.
\end{proof}

\begin{proof}[Proof of Proposition~\ref{prop:04-integer-bounds}]
For the estimates without a shift, take absolute values, use \(0\le A\le1\),
and apply Lemma~\ref{lem:04-quadratic} with \(F=\operatorname{Re}T\)
and infimum \(m_0\). The negative log sum is \(O(L)\) by
\eqref{eq:G-8} and the trace bound; the reference masses are bounded
above. The real trial estimate and \eqref{eq:G-5}, with
\eqref{eq:G-3}, prove the first two assertions.

On the full-shift range use the sharp version of \eqref{eq:G-18}, \eqref{eq:G-11}, and the infimum comparison to \(m_z\). Including the full Jacobian and both \(J_{\nu_\sigma}(w_0)\) gives the log bound there
\[
 -m_z+(7\gamma/8)L-\gamma L/2+o_\xi(L)+o(L)
\]
by \eqref{eq:G-9}. Taking \(\xi\) small, then applying the deterministic lower bounds on \(m_z\), and the separate negligible bad-region estimate (choose its power after the polynomial restrictions), proves \eqref{eq:G-16} and \eqref{eq:G-17}.
\end{proof}

